% Môn: Vật lý
% Lớp: 11
% Chương: Chương I. Dao động
% Bài: L11C1 Bài 1. Dao động điều hòa · Xác định li độ và pha của dao động tại một thời điểm xác định
% ===== Câu 74 =====
% ID: L11C1B1-148-TLN
% Mức: TH
\dangbt{Xác định li độ và pha của dao động tại một thời điểm xác định}
\begin{ex}
% ID: L11C1B1-148-TLN
% Mức: TH
Hai dao động cùng tần số có pha ban đầu $\pi/6$ và $\pi/2.$ Độ lệch pha bằng bao nhiêu lần $\pi$?
\shortans{01/03/2026}
\end{ex}

% ===== Câu 78 =====
% ID: L11C1B1-149-TLN
% Mức: VD
\begin{ex}
% ID: L11C1B1-149-TLN
% Mức: VD
Hai dao động cùng tần số có pha ban đầu $\pi/6$ và $\pi/2.$ Độ lệch pha bằng bao nhiêu lần $\pi$?
\shortans{01/03/2026}
\end{ex}

% ===== Câu 82 =====
% ID: L11C1B1-150-TLN
% Mức: VD
\begin{ex}
% ID: L11C1B1-150-TLN
% Mức: VD
Hai dao động cùng tần số có pha ban đầu $\pi/6$ và $\pi/2.$ Độ lệch pha bằng bao nhiêu lần $\pi$?
\shortans{01/03/2026}
\end{ex}

% ===== Câu 85 =====
% ID: L11C1B1-151-DS
% Mức: TH
\begin{ex}
% ID: L11C1B1-151-DS
% Mức: TH
Vật dao động theo $x=4\cos(2\pi t/1+\pi/3)$ cm. Xác định pha, li độ và chiều chuyển động tại $t=0$.
\choiceTF

\end{ex}

% ===== Câu 89 =====
% ID: L11C1B1-181-DS
% Mức: VD
\begin{ex}
% ID: L11C1B1-181-DS
% Mức: VD
Vật dao động theo $x=6\cos(2\pi t/2+\pi/3)$ cm. Xác định pha, li độ và chiều chuyển động tại $t=0$.
\choiceTF

\end{ex}

% ===== Câu 93 =====
% ID: L11C1B1-182-DS
% Mức: VD
\begin{ex}
% ID: L11C1B1-182-DS
% Mức: VD
Vật dao động theo $x=8\cos(2\pi t/4+\pi/3)$ cm. Xác định pha, li độ và chiều chuyển động tại $t=0$.
\choiceTF

\end{ex}

% ===== Câu 159 =====
% ID: L11C1B1-183-DS
% Mức: VD
\begin{ex}
% ID: L11C1B1-183-DS
% Mức: VD
Một vật dao động theo phương trình x = 4cos(2πt + π/3) cm. Hãy xác định tính đúng sai của các phát biểu sau.
\choiceTF
{\True Pha ban đầu của dao động là π/3 rad.}
{\True Tại t = 0, li độ của vật bằng $2\,\text{cm}$.}
{Tại t = 0, vật chuyển động theo chiều dương.}
{\True Tại t = 1/ $3\,\text{s}$, vật ở biên âm.}
\loigiai{
\begin{itemchoice}
\item (Đúng) Pha ban đầu của dao động là π/3 rad.
\item (Đúng) Tại t = 0, li độ của vật bằng $2\,\text{cm}$.
\item (Sai) Tại t = 0, vật chuyển động theo chiều dương.
\item (Đúng) Tại t = 1/ $3\,\text{s}$, vật ở biên âm.
\end{itemchoice}
}
\end{ex}

% ===== Câu 161 =====
% ID: L11C1B1-184-DS
% Mức: VD
\begin{ex}
% ID: L11C1B1-184-DS
% Mức: VD
Một vật dao động theo phương trình x = 4cos(2πt + π/3) cm. Hãy xác định tính đúng sai của các phát biểu sau.
\choiceTF
{\True Pha ban đầu của dao động là π/3 rad.}
{\True Tại t = 0, li độ của vật bằng $2\,\text{cm}$.}
{Tại t = 0, vật chuyển động theo chiều dương.}
{\True Tại t = 1/ $3\,\text{s}$, vật ở biên âm.}
\loigiai{
\begin{itemchoice}
\item (Đúng) Pha ban đầu của dao động là π/3 rad.
\item (Đúng) Tại t = 0, li độ của vật bằng $2\,\text{cm}$.
\item (Sai) Tại t = 0, vật chuyển động theo chiều dương.
\item (Đúng) Tại t = 1/ $3\,\text{s}$, vật ở biên âm.
\end{itemchoice}
}
\end{ex}

% ===== Câu 163 =====
% ID: L11C1B1-185-DS
% Mức: VD
\begin{ex}
% ID: L11C1B1-185-DS
% Mức: VD
Một vật dao động theo phương trình x = 4cos(2πt + π/3) cm. Hãy xác định tính đúng sai của các phát biểu sau.
\choiceTF
{\True Pha ban đầu của dao động là π/3 rad.}
{\True Tại t = 0, li độ của vật bằng $2\,\text{cm}$.}
{Tại t = 0, vật chuyển động theo chiều dương.}
{\True Tại t = 1/ $3\,\text{s}$, vật ở biên âm.}
\loigiai{
\begin{itemchoice}
\item (Đúng) Pha ban đầu của dao động là π/3 rad.
\item (Đúng) Tại t = 0, li độ của vật bằng $2\,\text{cm}$.
\item (Sai) Tại t = 0, vật chuyển động theo chiều dương.
\item (Đúng) Tại t = 1/ $3\,\text{s}$, vật ở biên âm.
\end{itemchoice}
}
\end{ex}

% ===== Câu 165 =====
% ID: L11C1B1-186-DS
% Mức: VD
\begin{ex}
% ID: L11C1B1-186-DS
% Mức: VD
Một vật dao động theo phương trình x = 6cos(πt - π/2) cm. Hãy xác định tính đúng sai của các phát biểu sau.
\choiceTF
{\True Pha ban đầu là -π/2 rad.}
{\True Tại t = 0, vật ở vị trí cân bằng.}
{\True Tại t = 0, vật chuyển động theo chiều dương.}
{\True Tại t = $0,5\,\text{s}$, vật ở biên dương.}
\loigiai{
\begin{itemchoice}
\item (Đúng) Pha ban đầu là -π/2 rad.
\item (Đúng) Tại t = 0, vật ở vị trí cân bằng.
\item (Đúng) Tại t = 0, vật chuyển động theo chiều dương.
\item (Đúng) Tại t = $0,5\,\text{s}$, vật ở biên dương.
\end{itemchoice}
}
\end{ex}

% ===== Câu 167 =====
% ID: L11C1B1-187-DS
% Mức: VD
\begin{ex}
% ID: L11C1B1-187-DS
% Mức: VD
Một vật dao động theo phương trình x = 6cos(πt - π/2) cm. Hãy xác định tính đúng sai của các phát biểu sau.
\choiceTF
{\True Pha ban đầu là -π/2 rad.}
{\True Tại t = 0, vật ở vị trí cân bằng.}
{\True Tại t = 0, vật chuyển động theo chiều dương.}
{\True Tại t = $0,5\,\text{s}$, vật ở biên dương.}
\loigiai{
\begin{itemchoice}
\item (Đúng) Pha ban đầu là -π/2 rad.
\item (Đúng) Tại t = 0, vật ở vị trí cân bằng.
\item (Đúng) Tại t = 0, vật chuyển động theo chiều dương.
\item (Đúng) Tại t = $0,5\,\text{s}$, vật ở biên dương.
\end{itemchoice}
}
\end{ex}

% ===== Câu 169 =====
% ID: L11C1B1-188-DS
% Mức: VD
\begin{ex}
% ID: L11C1B1-188-DS
% Mức: VD
Một vật dao động theo phương trình x = 6cos(πt - π/2) cm. Hãy xác định tính đúng sai của các phát biểu sau.
\choiceTF
{\True Pha ban đầu là -π/2 rad.}
{\True Tại t = 0, vật ở vị trí cân bằng.}
{\True Tại t = 0, vật chuyển động theo chiều dương.}
{\True Tại t = $0,5\,\text{s}$, vật ở biên dương.}
\loigiai{
\begin{itemchoice}
\item (Đúng) Pha ban đầu là -π/2 rad.
\item (Đúng) Tại t = 0, vật ở vị trí cân bằng.
\item (Đúng) Tại t = 0, vật chuyển động theo chiều dương.
\item (Đúng) Tại t = $0,5\,\text{s}$, vật ở biên dương.
\end{itemchoice}
}
\end{ex}

% ===== Câu 171 =====
% ID: L11C1B1-189-DS
% Mức: VD
\begin{ex}
% ID: L11C1B1-189-DS
% Mức: VD
Một vật dao động điều hòa được mô tả bằng phương trình li độ $x = 4\cos\left(2\pi t + \dfrac{\pi}{3}\right)$ (cm). Khảo sát các mốc thời gian chuyển động của vật:
\choiceTF
{\True Vật đạt li độ cực đại dương lần đầu tiên tại thời điểm $t = \dfrac{5}{6}$ s.}
{\True Vật đi qua vị trí cân bằng lần đầu tiên tại thời điểm $t = \dfrac{1}{12}$ s.}
{\True Khoảng thời gian ngắn nhất để vật di chuyển từ vị trí biên âm sang vị trí biên dương là $0,5$ s.}
{Tại thời điểm $t = 1$ s, vật đang ở vị trí cân bằng.}
\loigiai{
\begin{itemchoice}
\item (Đúng) Vật đạt li độ cực đại dương lần đầu tiên tại thời điểm $t = \dfrac{5}{6}$ s. (Vì): (Đúng) Vật đạt li độ cực đại dương lần đầu tiên tại thời điểm $t = \dfrac{5}{6}$ s
\item (Đúng) Vật đi qua vị trí cân bằng lần đầu tiên tại thời điểm $t = \dfrac{1}{12}$ s. (Vì): (Đúng) Vật đi qua vị trí cân bằng lần đầu tiên tại thời điểm $t = \dfrac{1}{12}$ s
\item (Đúng) Khoảng thời gian ngắn nhất để vật di chuyển từ vị trí biên âm sang vị trí biên dương là $0,5$ s. (Vì): (Đúng) Khoảng thời gian ngắn nhất để vật di chuyển từ vị trí biên âm sang vị trí biên dương là $0,5$ s
\item (Sai) Tại thời điểm $t = 1$ s, vật đang ở vị trí cân bằng. (Vì): o $T = \dfrac{2\pi}{2\pi} = 1$ s
\end{itemchoice}
}
\end{ex}

% ===== Câu 173 =====
% ID: L11C1B1-190-DS
% Mức: VD
\begin{ex}
% ID: L11C1B1-190-DS
% Mức: VD
Một vật dao động điều hòa được mô tả bằng phương trình li độ $x = 4\cos\left(2\pi t + \dfrac{\pi}{3}\right)$ (cm). Khảo sát các mốc thời gian chuyển động của vật:
\choiceTF
{\True Vật đạt li độ cực đại dương lần đầu tiên tại thời điểm $t = \dfrac{5}{6}$ s.}
{\True Vật đi qua vị trí cân bằng lần đầu tiên tại thời điểm $t = \dfrac{1}{12}$ s.}
{\True Khoảng thời gian ngắn nhất để vật di chuyển từ vị trí biên âm sang vị trí biên dương là $0,5$ s.}
{Tại thời điểm $t = 1$ s, vật đang ở vị trí cân bằng.}
\loigiai{
\begin{itemchoice}
\item (Đúng) Vật đạt li độ cực đại dương lần đầu tiên tại thời điểm $t = \dfrac{5}{6}$ s. (Vì): (Đúng) Vật đạt li độ cực đại dương lần đầu tiên tại thời điểm $t = \dfrac{5}{6}$ s
\item (Đúng) Vật đi qua vị trí cân bằng lần đầu tiên tại thời điểm $t = \dfrac{1}{12}$ s. (Vì): (Đúng) Vật đi qua vị trí cân bằng lần đầu tiên tại thời điểm $t = \dfrac{1}{12}$ s
\item (Đúng) Khoảng thời gian ngắn nhất để vật di chuyển từ vị trí biên âm sang vị trí biên dương là $0,5$ s. (Vì): (Đúng) Khoảng thời gian ngắn nhất để vật di chuyển từ vị trí biên âm sang vị trí biên dương là $0,5$ s
\item (Sai) Tại thời điểm $t = 1$ s, vật đang ở vị trí cân bằng. (Vì): o $T = \dfrac{2\pi}{2\pi} = 1$ s
\end{itemchoice}
}
\end{ex}

% ===== Câu 175 =====
% ID: L11C1B1-191-DS
% Mức: VD
\begin{ex}
% ID: L11C1B1-191-DS
% Mức: VD
Một vật dao động điều hòa được mô tả bằng phương trình li độ $x = 4\cos\left(2\pi t + \dfrac{\pi}{3}\right)$ (cm). Khảo sát các mốc thời gian chuyển động của vật:
\choiceTF
{\True Vật đạt li độ cực đại dương lần đầu tiên tại thời điểm $t = \dfrac{5}{6}$ s.}
{\True Vật đi qua vị trí cân bằng lần đầu tiên tại thời điểm $t = \dfrac{1}{12}$ s.}
{\True Khoảng thời gian ngắn nhất để vật di chuyển từ vị trí biên âm sang vị trí biên dương là $0,5$ s.}
{Tại thời điểm $t = 1$ s, vật đang ở vị trí cân bằng.}
\loigiai{
\begin{itemchoice}
\item (Đúng) Vật đạt li độ cực đại dương lần đầu tiên tại thời điểm $t = \dfrac{5}{6}$ s. (Vì): (Đúng) Vật đạt li độ cực đại dương lần đầu tiên tại thời điểm $t = \dfrac{5}{6}$ s
\item (Đúng) Vật đi qua vị trí cân bằng lần đầu tiên tại thời điểm $t = \dfrac{1}{12}$ s. (Vì): (Đúng) Vật đi qua vị trí cân bằng lần đầu tiên tại thời điểm $t = \dfrac{1}{12}$ s
\item (Đúng) Khoảng thời gian ngắn nhất để vật di chuyển từ vị trí biên âm sang vị trí biên dương là $0,5$ s. (Vì): (Đúng) Khoảng thời gian ngắn nhất để vật di chuyển từ vị trí biên âm sang vị trí biên dương là $0,5$ s
\item (Sai) Tại thời điểm $t = 1$ s, vật đang ở vị trí cân bằng. (Vì): o $T = \dfrac{2\pi}{2\pi} = 1$ s
\end{itemchoice}
}
\end{ex}

% ===== Câu 200 =====
% ID: L11C1B1-192-TN
% Mức: VD
\begin{ex}
% ID: L11C1B1-192-TN
% Mức: VD
Một vật dao động điều hòa với $x = 5\sin(4\pi t)$ cm. Tính li độ tại $t = \frac{1}{8}$ s.
\choice
{$x = 2{,}5$ cm}
{\True $x = 5$ cm}
{$x = -5$ cm}
{$x = 0$ cm}
\loigiai{
Thay $t = \frac{1}{8}$ s vào phương trình $x = 5\sin(4\pi t)$ cm:
  \begin{aligned}
 x &= 5 $\sin(4\pi\cdot\frac{1}{8})$ &= 5 $\sin(\frac{\pi}{2})$ &= 5\cdot 1 
&= 5 $\\text{cm}$.
 \end{aligned} 
 Vậy li độ tại $t = \frac{1}{8}$ s là $5\,\text{cm}$.
}
\end{ex}

% ===== Câu 201 =====
% ID: L11C1B1-193-TN
% Mức: VD
\begin{ex}
% ID: L11C1B1-193-TN
% Mức: VD
Một vật dao động điều hòa với $x = 5\sin(4\pi t)$ cm. Tính li độ tại $t = \frac{1}{8}$ s.
\choice
{$x = 2{,}5$ cm}
{\True $x = 5$ cm}
{$x = -5$ cm}
{$x = 0$ cm}
\loigiai{
Thay $t = \frac{1}{8}$ s vào phương trình $x = 5\sin(4\pi t)$ cm:
  \begin{aligned}
 x &= 5 $\sin(4\pi\cdot\frac{1}{8})$ &= 5 $\sin(\frac{\pi}{2})$ &= 5\cdot 1 
&= 5 $\\text{cm}$.
 \end{aligned} 
 Vậy li độ tại $t = \frac{1}{8}$ s là $5\,\text{cm}$.
}
\end{ex}

% ===== Câu 202 =====
% ID: L11C1B1-194-TN
% Mức: VD
\begin{ex}
% ID: L11C1B1-194-TN
% Mức: VD
Một vật dao động điều hòa với $x = 5\sin(4\pi t)$ cm. Tính li độ tại $t = \frac{1}{8}$ s.
\choice
{$x = 2{,}5$ cm}
{\True $x = 5$ cm}
{$x = -5$ cm}
{$x = 0$ cm}
\loigiai{
Thay $t = \frac{1}{8}$ s vào phương trình $x = 5\sin(4\pi t)$ cm:
  \begin{aligned}
 x &= 5 $\sin(4\pi\cdot\frac{1}{8})$ &= 5 $\sin(\frac{\pi}{2})$ &= 5\cdot 1 
&= 5 $\\text{cm}$.
 \end{aligned} 
 Vậy li độ tại $t = \frac{1}{8}$ s là $5\,\text{cm}$.
}
\end{ex}

% ===== Câu 203 =====
% ID: L11C1B1-195-TN
% Mức: VD
\begin{ex}
% ID: L11C1B1-195-TN
% Mức: VD
Dao động của một vật có phương trình $x = 4\cos(3\pi t - \frac{2\pi}{3})$ cm. Tại thời điểm $t = \frac{1}{3}$ s, li độ của vật là bao nhiêu?
\choice
{$x = -4$ cm}
{\True $x = 2$ cm}
{$x = 0$ cm}
{$x = -2$ cm}
\loigiai{
Thay $t = \frac{1}{3}$ s vào phương trình $x = 4\cos(3\pi t - \frac{2\pi}{3})$ cm:
  \begin{aligned}
 x &= 4 $\cos(3\pi\cdot\frac{1}{3}$ - $\frac{2\pi}{3})$ &= 4 $\cos(\pi$ - $\frac{2\pi}{3})$ &= 4 $\cos(\frac{\pi}{3})$ &= 4 $\cdot\frac{1}{2}$ &= 2 $\\text{cm}$.
 \end{aligned}
}
\end{ex}

% ===== Câu 204 =====
% ID: L11C1B1-196-TN
% Mức: VD
\begin{ex}
% ID: L11C1B1-196-TN
% Mức: VD
Dao động của một vật có phương trình $x = 4\cos(3\pi t - \frac{2\pi}{3})$ cm. Tại thời điểm $t = \frac{1}{3}$ s, li độ của vật là bao nhiêu?
\choice
{$x = -4$ cm}
{\True $x = 2$ cm}
{$x = 0$ cm}
{$x = -2$ cm}
\loigiai{
Thay $t = \frac{1}{3}$ s vào phương trình $x = 4\cos(3\pi t - \frac{2\pi}{3})$ cm:
  \begin{aligned}
 x &= 4 $\cos(3\pi\cdot\frac{1}{3}$ - $\frac{2\pi}{3})$ &= 4 $\cos(\pi$ - $\frac{2\pi}{3})$ &= 4 $\cos(\frac{\pi}{3})$ &= 4 $\cdot\frac{1}{2}$ &= 2 $\\text{cm}$.
 \end{aligned}
}
\end{ex}

% ===== Câu 205 =====
% ID: L11C1B1-197-TN
% Mức: VD
\begin{ex}
% ID: L11C1B1-197-TN
% Mức: VD
Dao động của một vật có phương trình $x = 4\cos(3\pi t - \frac{2\pi}{3})$ cm. Tại thời điểm $t = \frac{1}{3}$ s, li độ của vật là bao nhiêu?
\choice
{$x = -4$ cm}
{\True $x = 2$ cm}
{$x = 0$ cm}
{$x = -2$ cm}
\loigiai{
Thay $t = \frac{1}{3}$ s vào phương trình $x = 4\cos(3\pi t - \frac{2\pi}{3})$ cm:
  \begin{aligned}
 x &= 4 $\cos(3\pi\cdot\frac{1}{3}$ - $\frac{2\pi}{3})$ &= 4 $\cos(\pi$ - $\frac{2\pi}{3})$ &= 4 $\cos(\frac{\pi}{3})$ &= 4 $\cdot\frac{1}{2}$ &= 2 $\\text{cm}$.
 \end{aligned}
}
\end{ex}

% ===== Câu 206 =====
% ID: L11C1B1-198-TN
% Mức: VD
\begin{ex}
% ID: L11C1B1-198-TN
% Mức: VD
Phương trình dao động là $x = 6\cos(3\pi t - \frac{\pi}{4})$ cm. Tính vận tốc của vật tại thời điểm $t = \frac{1}{6}$ s.
\choice
{$v = 0$ cm/s}
{$v = 9\pi$ cm/s}
{\True $v = -9\pi$ cm/s}
{$v = 18\pi$ cm/s}
\loigiai{
Từ phương trình li độ $x = 6\cos(3\pi t - \frac{\pi}{4})$ cm, phương trình vận tốc là:
 $v = -6 \cdot 3\pi \sin(3\pi t - \frac{\pi}{4}) = -18\pi\sin(3\pi t - \frac{\pi}{4})\ \text{cm/s}.$
 Tại $t = \frac{1}{6}$ s:
  \begin{aligned}
 v &= -18 $\pi\sin(3\pi\cdot\frac{1}{6}$ - $\frac{\pi}{4})$ &= -18 $\pi\sin(\frac{\pi}{2}$ - $\frac{\pi}{4})$ &= -18 $\pi\sin(\frac{\pi}{4})$ &= -18 $\pi\cdot\frac{\sqrt{2}}{2}$ &\approx -12,73 $\\text{cm/s}$.
 \end{aligned} 
 Lưu ý: Số liệu hoặc đáp án của câu hỏi chưa thống nhất, cần kiểm tra lại.
}
\end{ex}

% ===== Câu 207 =====
% ID: L11C1B1-199-TN
% Mức: VD
\begin{ex}
% ID: L11C1B1-199-TN
% Mức: VD
Phương trình dao động là $x = 6\cos(3\pi t - \frac{\pi}{4})$ cm. Tính vận tốc của vật tại thời điểm $t = \frac{1}{6}$ s.
\choice
{$v = 0$ cm/s}
{$v = 9\pi$ cm/s}
{\True $v = -9\pi$ cm/s}
{$v = 18\pi$ cm/s}
\loigiai{
Từ phương trình li độ $x = 6\cos(3\pi t - \frac{\pi}{4})$ cm, phương trình vận tốc là:
 $v = -6 \cdot 3\pi \sin(3\pi t - \frac{\pi}{4}) = -18\pi\sin(3\pi t - \frac{\pi}{4})\ \text{cm/s}.$
 Tại $t = \frac{1}{6}$ s:
  \begin{aligned}
 v &= -18 $\pi\sin(3\pi\cdot\frac{1}{6}$ - $\frac{\pi}{4})$ &= -18 $\pi\sin(\frac{\pi}{2}$ - $\frac{\pi}{4})$ &= -18 $\pi\sin(\frac{\pi}{4})$ &= -18 $\pi\cdot\frac{\sqrt{2}}{2}$ &\approx -12,73 $\\text{cm/s}$.
 \end{aligned} 
 Lưu ý: Số liệu hoặc đáp án của câu hỏi chưa thống nhất, cần kiểm tra lại.
}
\end{ex}

% ===== Câu 208 =====
% ID: L11C1B1-200-TN
% Mức: VD
\begin{ex}
% ID: L11C1B1-200-TN
% Mức: VD
Phương trình dao động là $x = 6\cos(3\pi t - \frac{\pi}{4})$ cm. Tính vận tốc của vật tại thời điểm $t = \frac{1}{6}$ s.
\choice
{$v = 0$ cm/s}
{$v = 9\pi$ cm/s}
{\True $v = -9\pi$ cm/s}
{$v = 18\pi$ cm/s}
\loigiai{
Từ phương trình li độ $x = 6\cos(3\pi t - \frac{\pi}{4})$ cm, phương trình vận tốc là:
 $v = -6 \cdot 3\pi \sin(3\pi t - \frac{\pi}{4}) = -18\pi\sin(3\pi t - \frac{\pi}{4})\ \text{cm/s}.$
 Tại $t = \frac{1}{6}$ s:
  \begin{aligned}
 v &= -18 $\pi\sin(3\pi\cdot\frac{1}{6}$ - $\frac{\pi}{4})$ &= -18 $\pi\sin(\frac{\pi}{2}$ - $\frac{\pi}{4})$ &= -18 $\pi\sin(\frac{\pi}{4})$ &= -18 $\pi\cdot\frac{\sqrt{2}}{2}$ &\approx -12,73 $\\text{cm/s}$.
 \end{aligned} 
 Lưu ý: Số liệu hoặc đáp án của câu hỏi chưa thống nhất, cần kiểm tra lại.
}
\end{ex}

% ===== Câu 209 =====
% ID: L11C1B1-201-TN
% Mức: VD
\begin{ex}
% ID: L11C1B1-201-TN
% Mức: VD
Phương trình dao động của một vật là $x = 6\sin(\frac{\pi}{2}t)$ cm. Tính pha dao động tại thời điểm $t = 1$ s.
\choice
{\True $\varphi = \frac{\pi}{2}$ rad}
{$\varphi = 0$ rad}
{$\varphi = \pi$ rad}
{$\varphi = 2\pi$ rad}
\loigiai{
Pha dao động tại thời điểm $t$ được xác định bởi biểu thức trong hàm lượng giác.
 Với $t = 1$ s:
  \begin{aligned}
 $\varphi(1)$ &= $\frac{\pi}{2}\cdot$ 1 
&= $\frac{\pi}{2}\\text{rad}$.
 \end{aligned} 
 Vậy pha dao động tại $t = 1$ s là $\frac{\pi}{2}$ rad.
}
\end{ex}

% ===== Câu 210 =====
% ID: L11C1B1-202-TN
% Mức: VD
\begin{ex}
% ID: L11C1B1-202-TN
% Mức: VD
Phương trình dao động của một vật là $x = 6\sin(\frac{\pi}{2}t)$ cm. Tính pha dao động tại thời điểm $t = 1$ s.
\choice
{\True $\varphi = \frac{\pi}{2}$ rad}
{$\varphi = 0$ rad}
{$\varphi = \pi$ rad}
{$\varphi = 2\pi$ rad}
\loigiai{
Pha dao động tại thời điểm $t$ được xác định bởi biểu thức trong hàm lượng giác.
 Với $t = 1$ s:
  \begin{aligned}
 $\varphi(1)$ &= $\frac{\pi}{2}\cdot$ 1 
&= $\frac{\pi}{2}\\text{rad}$.
 \end{aligned} 
 Vậy pha dao động tại $t = 1$ s là $\frac{\pi}{2}$ rad.
}
\end{ex}

% ===== Câu 211 =====
% ID: L11C1B1-203-TN
% Mức: VD
\begin{ex}
% ID: L11C1B1-203-TN
% Mức: VD
Phương trình dao động của một vật là $x = 6\sin(\frac{\pi}{2}t)$ cm. Tính pha dao động tại thời điểm $t = 1$ s.
\choice
{\True $\varphi = \frac{\pi}{2}$ rad}
{$\varphi = 0$ rad}
{$\varphi = \pi$ rad}
{$\varphi = 2\pi$ rad}
\loigiai{
Pha dao động tại thời điểm $t$ được xác định bởi biểu thức trong hàm lượng giác.
 Với $t = 1$ s:
  \begin{aligned}
 $\varphi(1)$ &= $\frac{\pi}{2}\cdot$ 1 
&= $\frac{\pi}{2}\\text{rad}$.
 \end{aligned} 
 Vậy pha dao động tại $t = 1$ s là $\frac{\pi}{2}$ rad.
}
\end{ex}

% ===== Câu 212 =====
% ID: L11C1B1-204-TN
% Mức: VD
\begin{ex}
% ID: L11C1B1-204-TN
% Mức: VD
Phương trình dao động là $x = 3\cos(2t + \frac{\pi}{4})$ cm. Tính pha dao động tại thời điểm $t = 2$ s.
\choice
{$\varphi = \frac{\pi}{4}$ rad}
{$\varphi = 2$ rad}
{\True $\varphi = 4,25$ rad}
{$\varphi = \pi$ rad}
\loigiai{
Pha dao động tại thời điểm $t = 2$ s là:
  \begin{aligned}
 $\varphi(2)$ &= 2 \cdot 2 + $\frac{\pi}{4}$ &= 4 + $\frac{\pi}{4}$ &\approx 4 + 0,785 
&\approx 4,785 $\\text{rad}$.
 \end{aligned} 
 Lưu ý: Số liệu hoặc đáp án của câu hỏi chưa thống nhất, cần kiểm tra lại.
}
\end{ex}

% ===== Câu 213 =====
% ID: L11C1B1-205-TN
% Mức: VD
\begin{ex}
% ID: L11C1B1-205-TN
% Mức: VD
Phương trình dao động là $x = 3\cos(2t + \frac{\pi}{4})$ cm. Tính pha dao động tại thời điểm $t = 2$ s.
\choice
{$\varphi = \frac{\pi}{4}$ rad}
{$\varphi = 2$ rad}
{\True $\varphi = 4,25$ rad}
{$\varphi = \pi$ rad}
\loigiai{
Pha dao động tại thời điểm $t = 2$ s là:
  \begin{aligned}
 $\varphi(2)$ &= 2 \cdot 2 + $\frac{\pi}{4}$ &= 4 + $\frac{\pi}{4}$ &\approx 4 + 0,785 
&\approx 4,785 $\\text{rad}$.
 \end{aligned} 
 Lưu ý: Số liệu hoặc đáp án của câu hỏi chưa thống nhất, cần kiểm tra lại.
}
\end{ex}

% ===== Câu 214 =====
% ID: L11C1B1-206-TN
% Mức: VD
\begin{ex}
% ID: L11C1B1-206-TN
% Mức: VD
Phương trình dao động là $x = 3\cos(2t + \frac{\pi}{4})$ cm. Tính pha dao động tại thời điểm $t = 2$ s.
\choice
{$\varphi = \frac{\pi}{4}$ rad}
{$\varphi = 2$ rad}
{\True $\varphi = 4,25$ rad}
{$\varphi = \pi$ rad}
\loigiai{
Pha dao động tại thời điểm $t = 2$ s là:
  \begin{aligned}
 $\varphi(2)$ &= 2 \cdot 2 + $\frac{\pi}{4}$ &= 4 + $\frac{\pi}{4}$ &\approx 4 + 0,785 
&\approx 4,785 $\\text{rad}$.
 \end{aligned} 
 Lưu ý: Số liệu hoặc đáp án của câu hỏi chưa thống nhất, cần kiểm tra lại.
}
\end{ex}

% ===== Câu 215 =====
% ID: L11C1B1-207-TN
% Mức: VD
\begin{ex}
% ID: L11C1B1-207-TN
% Mức: VD
Một dao động điều hòa có phương trình $x = 8\sin(\pi t)$ cm. Tính gia tốc của vật tại vị trí li độ $x = 4$ cm.
\choice
{$a = -0,5\pi^2$ cm/s $^2$}
{$a = -\pi^2$ cm/s $^2$}
{\True $a = -\frac{\pi^2}{2}$ cm/s $^2$}
{$a = -2\pi^2$ cm/s $^2$}
\loigiai{
Từ phương trình $x = 8\sin(\pi t)$ cm, ta có $A = 8$ cm và $\omega = \pi$ rad/s.
 Gia tốc trong dao động điều hòa là $a = -\omega^2 x$.
 Tại $x = 4$ cm:
  \begin{aligned}
 a &= - $\omega^2$ x 
&= - $\pi^2\cdot$ 4 
&= -4 $\pi^2\\text{cm/s}^2$.
 \end{aligned} 
 Lưu ý: Số liệu hoặc đáp án của câu hỏi chưa thống nhất, cần kiểm tra lại.
}
\end{ex}

% ===== Câu 216 =====
% ID: L11C1B1-208-TN
% Mức: VD
\begin{ex}
% ID: L11C1B1-208-TN
% Mức: VD
Một dao động điều hòa có phương trình $x = 8\sin(\pi t)$ cm. Tính gia tốc của vật tại vị trí li độ $x = 4$ cm.
\choice
{$a = -0,5\pi^2$ cm/s $^2$}
{$a = -\pi^2$ cm/s $^2$}
{\True $a = -\frac{\pi^2}{2}$ cm/s $^2$}
{$a = -2\pi^2$ cm/s $^2$}
\loigiai{
Từ phương trình $x = 8\sin(\pi t)$ cm, ta có $A = 8$ cm và $\omega = \pi$ rad/s.
 Gia tốc trong dao động điều hòa là $a = -\omega^2 x$.
 Tại $x = 4$ cm:
  \begin{aligned}
 a &= - $\omega^2$ x 
&= - $\pi^2\cdot$ 4 
&= -4 $\pi^2\\text{cm/s}^2$.
 \end{aligned} 
 Lưu ý: Số liệu hoặc đáp án của câu hỏi chưa thống nhất, cần kiểm tra lại.
}
\end{ex}

% ===== Câu 217 =====
% ID: L11C1B1-209-TN
% Mức: VD
\begin{ex}
% ID: L11C1B1-209-TN
% Mức: VD
Một dao động điều hòa có phương trình $x = 8\sin(\pi t)$ cm. Tính gia tốc của vật tại vị trí li độ $x = 4$ cm.
\choice
{$a = -0,5\pi^2$ cm/s $^2$}
{$a = -\pi^2$ cm/s $^2$}
{\True $a = -\frac{\pi^2}{2}$ cm/s $^2$}
{$a = -2\pi^2$ cm/s $^2$}
\loigiai{
Từ phương trình $x = 8\sin(\pi t)$ cm, ta có $A = 8$ cm và $\omega = \pi$ rad/s.
 Gia tốc trong dao động điều hòa là $a = -\omega^2 x$.
 Tại $x = 4$ cm:
  \begin{aligned}
 a &= - $\omega^2$ x 
&= - $\pi^2\cdot$ 4 
&= -4 $\pi^2\\text{cm/s}^2$.
 \end{aligned} 
 Lưu ý: Số liệu hoặc đáp án của câu hỏi chưa thống nhất, cần kiểm tra lại.
}
\end{ex}
